\everymath{\displaystyle}

%%%%%%%%%%%%%%%%%%%%%%%%%%%%%%%%%%%%%%%%
%           main body
%%%%%%%%%%%%%%%%%%%%%%%%%%%%%%%%%%%%%%%%

%-----------------------------------
%	TITLE PAGE
%-----------------------------------

\title[第六部分\\
Nakayama 引理]{\texorpdfstring{\LARGE \bfseries 第六部分\quad Nakayama 引理}{第六部分 Nakayama 引理}}
\author[抽象代数 2]{\Large \bfseries 抽象代数 2}
\date{}

%------------------------------------------------

\begin{document}

%------------------------------------------------

\begin{frame}

\titlepage % Print the title page as the first slide

\end{frame}

%------------------------------------------------

\section[定理]{Nakayama 引理}

%------------------------------------------------

\begin{frame}{Nakayama 引理 (定理 2)}

\begin{thm}[Nakayama 引理]
$R$ 为幺环, $I \subset J(R)$ 为理想, $M$ 为有限生成的 $R$-模, $I \cdot M$
为 $I$ 在 $M$ 中生成的子模, $M/I \cdot M$ 为商模, 则:
\begin{itemize}
\item[(i)] $I \cdot M = M \iff M = 0$;
\item[(ii)] 任取 $H \subset M$ 为子模, $I \cdot M + H = M \iff H = M$;
\item[(iii)] $M = \langle \alpha_1, \alpha_2, \dots, \alpha_n \rangle
  \iff M/I \cdot M = \langle \bar{\alpha}_1, \bar{\alpha}_2, \dots,
  \bar{\alpha}_n \rangle$ 作为 $R/I$ 模.
\end{itemize}
\end{thm}

\end{frame}

%------------------------------------------------

\begin{frame}{证明的准备 (1) (2)}

\startproof[bg=yellow, fg=red](证明)
(1). 任一幺环 $R$ 中总存在至少一个极大左理想 $I$,
(思路) 即单模 $M = R/I$ 总是存在的.

\pause

(2). 任一 $R$ 的左理想都包含在 $R$ 的一个极大左理想中:
令 $N$ 为 $R$ 的一个左理想, 那么 $N$ 为 $R$ 的子模 ($R$ 作为 $R$ 模)
\[
\pi: N \to R/N \quad \text{为 } R\text{-自然模同态},
\]
$R$ 中包含 $\ker \pi = N$ 的子模与 $R/N$ 全部子模是 $1$-$1$ 对应的.
若 $R/N$ 有极大 $R$-子模 $\bar{H}$ 存在,
则 $N \subset \pi^{-1}(\bar{H}) \subset R$, $\pi^{-1}(\bar{H})$ 为极大子模.
\[
\begin{cases}
R = \langle 1 \rangle \text{ 为循环模}, \\
R/N = \langle \bar{1} \rangle \text{ 为循环模}.
\end{cases}
\]

\end{frame}

%------------------------------------------------

\begin{frame}{证明的准备 (3) (4)}

(3). 任一 $R$-循环模 $M$, $M \neq 0$ $\Longrightarrow$ $M$ 有极大子模:
令 $M = \langle \alpha \rangle$, $\alpha \neq 0$,
令 $S = \{ H \mid H \subset M \text{ 为子模},\ \alpha \notin H \}$,
可以由 Zorn 引理验证 $S$ 中有极大元素, 为 $M$ 的一个极大子模.

\pause

(4). $a \in R$, 为左不可逆元, 则 $a$ 含于 $R$ 的一个极大左理想中:
$R$ 为 $R$-模, $a \in R$, $\langle a \rangle$ 为子模, 也是一个左理想,
其中 $\langle a \rangle = Ra = \{ xa \mid x \in R \}$.
$a$ 左不可逆 $\Longrightarrow \langle a \rangle \neq R$,
由 (3) 知 $\langle a \rangle$ 包含于 $R$ 的一个极大左理想中.

\end{frame}

%------------------------------------------------

\section[根的元素刻画]{根的元素刻画}

%------------------------------------------------

\begin{frame}{根的元素刻画}

\begin{lem}[根的元素刻画]
$R$ 为幺环, $J(R) = \bigcap\limits_{I \text{ 为极大左理想}} I$,
则
\[
J(R) = \big\{ x \in R \;\big|\; \forall\ a \in R,\ 1 - ax \text{ 左可逆} \big\} .
\]
\end{lem}

\end{frame}

%------------------------------------------------

\begin{frame}{根的元素刻画的证明}

\startproof[bg=yellow, fg=red](证明)
任取 $x \in R$, 若 $\exists\ a \in R$, 使得 $1 - ax$ 左不可逆,
则存在极大左理想 $N$, 使得 $1 - ax \in N$, 由 $x \in N$ 知 $ax \in N$,
那么 $1 \in N$, 那么 $N = R$, 矛盾.
\pause
反过来, 若对 $x \in R$, $\forall\ a \in R$, $1 - ax$ 左可逆,
假设存在极大左理想 $N$, 使得 $x \notin N$, 那么 $N + \langle x \rangle = R$,
那么 $\exists\ a \in R$, 使得 $y \in N$, 使得 $y + ax = 1$.
$y = 1 - ax$ 左可逆, 即知 $N = R$, 矛盾.
综上知
$J(R) = \{ x \in R \mid \forall\ a \in R,\ 1 - ax \text{ 左可逆} \}$.
\qed

\end{frame}

%------------------------------------------------

\section[Nakayama 的证明]{Nakayama 引理的证明}

%------------------------------------------------

\begin{frame}{Nakayama 引理的证明 (1): (i)}

\startproof[bg=yellow, fg=red](Nakayama 引理 (定理 2) 的证明)
(1) $M = \langle \alpha_1, \alpha_2, \dots, \alpha_m \rangle$,
令 $m$ 为使上式成立的最小正整数, $m \neq 0$.
任取 $I \subset J(R)$ 为 $R$ 的理想, 假设 $I \cdot M = M$.
\[
M = \Big\{ \sum_i a_i \cdot \alpha_i \;\Big|\; a_i \in R \Big\},
\qquad
I \cdot M = \Big\{ \sum_i a_i \cdot \alpha_i \;\Big|\; a_i \in I \Big\},
\]
那么 $\alpha_m = \sum_i a_i \cdot \alpha_i$, 其中 $a_i \in I$,
那么 $(1 - a_m)\alpha_m = \sum_{i=1}^{m-1} a_i \cdot \alpha_i$,
由根的元素刻画引理知 $1 - a_m$ 左可逆,
\pause
即知 $\alpha_m = \sum_{i=1}^{m-1} b_i \cdot \alpha_i$, $b_i \in I$,
那么 $M = \langle \alpha_1, \dots, \alpha_{m-1} \rangle$ 与 $m$
的极小性矛盾.
\qed

\end{frame}

%------------------------------------------------

\begin{frame}{Nakayama 引理的证明 (2): (ii)}

\startproof[bg=yellow, fg=red](Nakayama 引理 (定理 2) 的证明, 续)
(ii) 任取 $M$ 的子模 $H$, $M/H$ 也是有限生成的 $R$-模,
下证 $I \cdot (M/H) = (I \cdot M + H)/H$:
$\pi: M \to M/H$ 为自然同态, $\pi|_{I \cdot M}: I \cdot M \to I \cdot (M/H)$
为满的模同态.
由同态基本定理,
$I \cdot M / \ker(\pi|_{I \cdot M}) \cong I \cdot (M/H)$,
而 $\ker(\pi|_{I \cdot M}) = I \cdot M \cap H$,
$I \cdot M / (I \cdot M \cap H) \cong (I \cdot M + H)/H$.
\pause
由 (1) 知
\[
I \cdot (M/H) = M/H \iff IM + H = M \iff M/H = 0
\iff H = M,
\]
即 $I \cdot M + H = M \iff H = M$.
\qed

\end{frame}

%------------------------------------------------

\begin{frame}{Nakayama 引理的证明 (3): (iii)}

\startproof[bg=yellow, fg=red](Nakayama 引理 (定理 2) 的证明, 续)
(iii) $\{ \alpha_1, \alpha_2, \dots, \alpha_n \}$ 为 $M$ 的一组生成元,
$M \xrightarrow{\ \pi\ } M/I \cdot M$, $\forall\ \alpha \in M$,
$\alpha = \sum_i a_i \cdot \alpha_i$, $a_i \in R$,
那么 $\bar{\alpha} = \pi(\alpha) = \sum_i \overline{a_i \alpha_i}
= \sum_i a_i \bar{\alpha}_i$
$\Longrightarrow \{ \bar{\alpha}_1, \dots, \bar{\alpha}_n \}$ 为
$M/I \cdot M$ 的一组生成元.
反过来, 若 $\{ \bar{\alpha}_1, \dots, \bar{\alpha}_n \}$ 为
$M/I \cdot M$ 的一组生成元 (作为 $R/I$-模),
那么 $\{ \bar{\alpha}_1, \dots, \bar{\alpha}_n \}$ 为 $M/I \cdot M$
作为 $R$-模的一组生成元.
\pause
$M \xrightarrow{\ \pi\ } M/I \cdot M$ 为自然同态, $\ker \pi = I \cdot M$,
令 $H = \langle \alpha_1, \dots, \alpha_n \rangle \subset M$,
由于 $\pi^{-1}(M/I \cdot M) = M$, $\pi(H) = M/I \cdot M$,
那么 $\pi^{-1}(M/I \cdot M) = H + I \cdot M$, 即 $H + I \cdot M = M$,
由 (ii) 知 $M = H = \langle \alpha_1, \dots, \alpha_n \rangle$.
\qed

\end{frame}

%------------------------------------------------

\end{document}
